These studies test sensitivity and retain adverse results. They are not the central novelty claim.

\subsection{Main real-data comparison}\label{main-real-data-comparison}

Table~\ref{tab:broad} gives mean fractions of pair-items evaluated at \(\alpha=.05\) and \(\epsilon=.02\). Full census is 1.000 for every matrix. The betting column is the direct-betting comparator; all other displayed methods belong to the core protocol.

\begin{table}[H]
\centering
\caption{Mean evaluated fractions for the main public cohort.}\label{tab:broad}
\begin{tabular}{lrrrrr}
\toprule
Matrix & EB grid & \shortstack{Predictable\\EB} & \shortstack{Hoeffding\\grid} & \shortstack{Serfling\\spending} & \shortstack{Direct\\betting}\\\midrule
ARC & .575 & .528 & .727 & .793 & .481\\
Composite & .294 & .326 & .398 & .447 & .261\\
GSM8K & .365 & .346 & .439 & .488 & .310\\
HellaSwag & .356 & .346 & .603 & .669 & .291\\
MMLU-Pro & .149 & .160 & .227 & .261 & .129\\
TruthfulQA & .491 & .443 & .610 & .671 & .411\\
WinoGrande & .685 & .652 & .844 & .902 & .588\\\bottomrule
\end{tabular}
\end{table}

\begin{figure}[H]
\centering
\includegraphics[width=.95\textwidth]{figures/real_cost.png}
\caption{Certified-method cost comparison across the seven public matrices.}\label{fig:random}
\end{figure}

The EB-grid macro-average is .41635. It produces no observed errors across the 21,000 main-margin replays. Direct betting is cheaper in each matrix and produces 11 errors across the same 21,000 replays; this remains an empirical statement, not an assertion that its true risk is that pooled frequency. Its proven guarantee is .05 per frozen vector. Neither the grid nor predictable EB dominates the other across matrices. Therefore the EB grid should not be promoted as a superior new algorithm.

Equivalence-specific costs are much less favorable than the aggregate. For EB-grid, the mean evaluated fractions among equivalent pairs are .978 on ARC, .980 on GSM8K, .983 on TruthfulQA, and .960 on WinoGrande. They are lower on HellaSwag (.632), MMLU-Pro (.637), and the composite (.791). Large savings over a random pair cohort can conceal expensive equivalence decisions. No claim about close leading models follows from randomly selected pairs alone.

\subsection{Pairing and stratification}\label{pairing-and-stratification}

Separate-arm confidence sequences are more expensive than paired EB-grid in all seven main-margin matrices. This is an expected benefit of exploiting model correlation, already central to SySRs; it is not a new finding about pairing in general.

On the exploratory public-data cohort, simple combined-stratum intervals usually cost more than an unstratified EB-grid. For the composite, the mean ratio of stopping costs with 38 task strata is about 21.5. This statistic is the average of per-replay ratios, not the ratio of mean costs, and is amplified by very early unstratified decisions. It should not be interpreted as a universal multiplicative penalty.

Controlled simulations show that stratification can nevertheless help. With two opposite constant strata, it evaluates .341 of the population versus .982 without strata. With a large constant stratum and a small noisy stratum, estimated-Neyman allocation evaluates .400 versus .772 unstratified and .859 under proportional allocation. With homogeneous mixed strata, it is worse (.964 versus .898). These examples establish a contingent trade-off between within-stratum variance reduction, allocation, and simultaneous-inference overhead. They do not establish optimal adaptive allocation.

\subsection{Stronger joint stratified inference}\label{stronger-joint-stratified-inference}

The cost of combining separate stratum intervals is not intrinsic to stratification. We additionally implement joint union-intersection exponential tests, using the same fixed bet grid but optimizing the nuisance vector of stratum means over the global null and deterministic completion box. The optimization reduces to fractional knapsack. This is a specialization of established stratified inference, with a proof in Appendix~\ref{theory:9}; we do not claim a new general testing principle. Both proportional and estimated-Neyman allocations are predictable. Inspections occur every 16 items, with all evaluations charged.

On the same 20-pair, five-permutation cohort, 5,800 joint-stratified replays give no observed errors. For the composite's 38 task strata, mean evaluated fractions are .186 with proportional allocation and .545 with estimated-Neyman allocation, compared with .764 and .752 for the respective separate-interval procedures. The unstratified EB-grid fraction on this cohort is .172. Thus the large separate-interval penalty mostly disappears with a stronger inference construction, but these task strata still do not establish a savings advantage over the uniform comparator. Across all seven matrices and each tested stratum definition, proportional allocation is cheaper than estimated-Neyman allocation for this joint test. This is evidence about these allocations and test statistics, not an impossibility result for adaptive sampling. Independent linear-programming comparisons and exhaustive small-population checks supplement the mathematical proof.

Matched controls use the same pairs and seeds: a one-stratum version of the joint test with the same 16-item inspection schedule, and direct uniform-pair betting. They add 1,400 replays. The one-stratum joint mean fraction on the composite is .169, versus .186 for task-stratified proportional inference. Direct betting is .146 on that cohort and is cheaper than every tested joint-stratified configuration in all seven matrices. Consequently, the observed repair of the separate-interval penalty is not evidence that stratification beats a strong uniform comparator. The direct test inspects every item, so this last comparison includes different inspection schedules; the matched one-stratum joint control isolates that scheduling difference.

\subsection{Prediction, close pairs, and heterogeneous costs}\label{app:prediction}\label{prediction-close-pairs-and-heterogeneous-costs}

We evaluate a CELEUS finite-pool signal adaptation on 30 preselected pairs per matrix and ten permutations, using 64 charged pilot pair-items to select nearest historical response profiles from 32 disjoint reference models. Uniform constant prediction, uniform historical prediction, and variance-weighted historical prediction share the same signal-based testing code. All sampling probabilities are accounted for in the signal; predictor accuracy is not needed for validity. This is not a reproduction of the authors' full predictor configuration. Across 6,300 replays, historical profiles usually increase cost. On HellaSwag, variance-weighted prediction reduces the mean fraction from .272 to .262; on ARC it increases it from .517 to .640. We report this negative predictor result without inferring that prediction assistance generally fails. Pilot and reference-model provenance are retained in the repository.

Random pairs are supplemented with two retrospectively selected cohorts: adjacent pairs among the top 21 released model rows, and the 20 pairs closest to the practical margin in absolute full-matrix mean difference. Each receives 100 permutations and two methods, giving 56,000 replays. Selection uses full-matrix truth to define the diagnostic cohort; no unrevealed target responses enter the algorithm. These results describe selected finite pairs, not a prospective model-selection procedure. Margin-nearest pairs require mean fractions .9974--.9999 for EB-grid and .9766--.9974 for direct betting across matrices. Top-adjacent costs vary much more: direct betting ranges from .157 on HellaSwag to .910 on GSM8K. The largest dataset/cohort direct-betting error frequency is .0135, below its nominal .05, although pooled frequencies do not certify individual-pair risk. A separate 69,000-replay direct-betting simulation audit has a maximum cell error frequency .020; per-cell intervals are provided.

Finally, a controlled two-stratum cost study charges a tenfold cost difference, with 300 repetitions for each method and orientation of the cost assignment. In the large-constant/small-noisy setting, estimated-Neyman allocation's mean item fraction is .400, while its cost fraction is .684 when the noisy stratum is expensive and .340 when the constant stratum is expensive. Uniform sampling costs approximately .768 and .773 respectively. These are synthetic cost units; the public matrices lack token and dollar logs. The comparison demonstrates why item savings should not be reported as monetary savings.

\subsection{Historical-only low-rank prediction sensitivity}\label{historical-only-low-rank-prediction-sensitivity}

To test whether the nearest-profile predictor explains the unfavorable prediction results, we also evaluate a squared-loss low-rank surrogate. The basis uses 512 historical model rows; rank and ridge penalty are selected on 64 other historical rows using prediction error outside a simulated 64-item pilot. All historical IDs are disjoint from target IDs. Candidate ranks are 0, 4, 8, 16 and 32; penalties are .1, 1, 10 and 100. Each target then pays for its own 64-item pilot. The predictor and residual-based weights remain frozen. This is an explicitly different implementation from PULSE's logistic factorization, compared under our common finite-pool inference layer.

The 4,200 low-rank replays use the same target pairs and pilot seeds as Section~\ref{app:prediction} and produce no observed errors. Selected ranks range from 4 to 16. Uniform low-rank evaluated fractions are .566 (ARC), .139 (composite), .304 (GSM8K), .282 (HellaSwag), .139 (MMLU-Pro), .629 (TruthfulQA), and .621 (WinoGrande). They improve on the uniform nearest-profile predictor in all seven matrices, but improve on the uniform constant-prediction control only slightly on MMLU-Pro (.139 versus .141). Residual-weighted sampling is more expensive than uniform low-rank sampling in every matrix. Thus this stronger predictor does not reverse the cost conclusion for the implemented signal test. Figure~\ref{fig:followup} reports paired differences and descriptive pair-cluster bootstrap intervals; neither this cohort nor its intervals establishes universal performance. Held-out target prediction errors are recorded strictly as post-run diagnostics, never used for selection or stopping.

\begin{figure}[H]
\centering
\includegraphics[width=.95\textwidth]{figures/followup_costs.png}
\caption{Panel A compares direct betting on two retrospective cohorts (20 pairs and 100 permutations per cohort and matrix). Panel B compares inference and allocation on the exploratory cohort (20 pairs and five permutations); the two joint methods inspect every 16 items, whereas direct betting inspects every item. Panel C shows low-rank minus constant-predictor cost on 30 matched pairs with ten permutations each. Error bars in C are descriptive 95\% percentile bootstrap intervals resampling pairs as clusters, using 5,000 draws. Positive differences mean higher cost. Panels A and B show finite-cohort means without uncertainty intervals. Cohorts differ across panels and must not be pooled as independent model draws.}\label{fig:followup}
\end{figure}


